\subsection{Ambiente estudado}

O estudo foi realizado em \textcolor{red}{\textbf{descrever a residencia: numero de comodos, area aproximada em m$^2$, tipo de construcao (alvenaria, madeira etc.) e  ateriais relevantes, como armarios metalicos ou paredes com tubulacao}}. A Figura~\ref{fig:planta} apresenta a planta baixa do ambiente, com a posicao original do ponto de acesso indicada.

\begin{figure}[ht]
  \centering
  % \includegraphics[width=.8\textwidth]{planta_baixa.png}
 \textcolor{red}{ \textbf{inserir aqui a planta baixa da residencia, com a posicao
  original do AP marcada }}
  \caption{Planta baixa do ambiente estudado, com a posicao original do ponto de acesso.}
  \label{fig:planta}
\end{figure}

O ponto de acesso utilizado foi \textcolor{red}{\textbf{modelo do roteador/AP, frequencia (2,4~GHz e/ou 5~GHz), ganho nominal da antena }}. As medicoes de potencia recebida foram feitas com \textcolor{red}{\textbf{nome do aplicativo ou ferramenta usada}}.

\subsection{Malha de medicao}

Foi definida uma malha de pontos com espacamento regular de \textcolor{red}{\textbf{valor em metros }} entre pontos adjacentes, adensada junto a paredes, portas e ao proprio ponto de acesso. Cada ponto foi numerado e sua posicao registrada em coordenadas $(x, y)$ relativas a uma origem escolhida na planta baixa, de modo que a mesma malha pudesse ser repetida de forma identica apos o reposicionamento do ponto de acesso. 
\subsection{Protocolo de coleta}

Em cada ponto da malha foram coletadas \textcolor{red}{\textbf{numero de amostras}} leituras de potencia recebida (dBm), ao longo de alguns segundos, das quais foi registrada a mediana. Foram padronizados a altura do aparelho de medicao, sua orientacao e a pessoa responsavel pela coleta, de modo que essas variaveis nao interferissem no resultado. O horario e as condicoes do ambiente (presenca de pessoas, portas abertas ou fechadas) foram registrados para cada campanha de medicao, garantindo que a segunda campanha fosse realizada em condicoes comparaveis a primeira.
\subsection{Etapa de medicao}

A primeira etapa foi realizada com o ponto de acesso em sua posicao original, produzindo o mapa de calor da Secao~\ref{subsec:heatMap}. Com base no padrao de radiacao esperado e nos obstaculos identificados no ambiente (paredes, moveis metalicos, distancia aos comodos mais afastados), foi formulada a hipotese de que reposicionar o ponto de acesso para \textcolor{red}{\textbf{descrever a nova posicao proposta e o raciocinio, o que essa posicao melhora e o que ela piora}} produziria uma cobertura mais equilibrada. O ponto de acesso foi entao movido para essa nova posicao e uma segunda campanha de medicao foi realizada, repetindo a mesma malha, o mesmo aparelho e o mesmo protocolo da primeira.